\section*{Discussion}

Two statements about $\Fovs$ have been run together, and they need to be
separated. The first is that the Atlantic freshwater budget at 34.5\degS{} sits
on the side where the salt-advection feedback is self-reinforcing. Three
reanalyses spanning two ocean models and three assimilation methods agree on
that, with intervals clear of zero at every block length we tried, at every
latitude from 28 to 34.5\degS{} in the two products run on the latitude ladder, and in the most recent decade of each record.
The second is that the observed decline in $\Fovs$ measures a weakening
overturning. That statement does not survive contact with the same data.

The reason is a rate comparison rather than a ratio. In the two products that
decline, $\Fovs$ is changing at 36\% and 52\% of its own magnitude per decade
while the volume transport that carries it changes at 2.6\% and 0.6\%. Passed
through equation~\eqref{eq:factorisation} this makes the transport contribution
between 1\% and 18\% of the trend depending on how the overturning is defined
and over which window, with individual intervals reaching 27\%, and in most
configurations it has the opposite sign to the total. It is a minority
contribution under every test we ran in which the share is identified, though it
is not definition-independent:
it rises systematically as the definition of the overturning narrows from the
whole northward limb to the streamfunction maximum. Studies that read a falling
$\Fovs$ as a proxy for AMOC weakening, including the early-warning arguments
built on the $\Fovs$ minimum\cite{vanWesten2024} and an earlier, unpublished
version of our own analysis of these data, have attributed to the circulation a
change that the circulation did not make.

What did make it is a widening of the salinity contrast between the limbs,
carried by the northward limb while the southward limb does not change. How much
of that is a water-mass change rather than a redistribution of transport in
depth differs between products: in GLORYS12V1 it is 94\%, in ORAS5 only about
two thirds, with the salinity-only part there not individually significant. The
coarse upper-to-abyssal redistribution is small in both, and excluding the
abyssal cell makes the salinification larger rather than smaller, so what
remains in ORAS5 is a rearrangement of transport within the upper limb.

The part of this that observations can check, they support. Holding the
circulation fixed and letting only the EN4 objective analysis vary, the
northward limb salinifies at $+0.026$ to $+0.028$~PSU per decade
($p\le0.0005$) while the southward limb is flat. The limb asymmetry is
therefore present in a salinity field that no model integrates forward, and the
same construction returns a negative $\Fovs$, $-0.080$ to $-0.128$~Sv, from
observed salinity. The rate is less secure: ORAS5 runs about twice the
observation-only value over the same years while GLORYS12V1 is within 30\% of
it, and part of that gap is expected because EN4 relaxes to climatology where
profiles are sparse. The direction should be trusted; the ORAS5 rate should
not.
It also reframes what the diagnostic is sensing. A salinifying upper limb widens
the limb contrast and drives $\Fovs$ downward regardless of what the overturning
does, which is why $\Fovs$ and the transport are so weakly coupled on the
timescales the records resolve. It is consistent in direction with the South
Atlantic salinity accumulation attributed to a weakening
overturning\cite{Zhu2020,Zhu2023}, but our decomposition locates the change in
the transport-weighted upper limb rather than in the surface field, and we find
no transport change to accompany it.

Two things temper this. The first is that $\Fovs$ is not the freshwater budget
of the basin. The azonal component at this section is several times larger, of
the opposite sign, and trending upward by more than $\Fovs$ is trending
downward: significantly so in ORAS5, and marginally in GLORYS12V1 ($p=0.049$ on
8.5 effective degrees of freedom). Whether the basin as a whole is
moving deeper into the self-reinforcing regime is therefore not settled by
$\Fovs$ alone, and we make no claim that it is. Nor do we evaluate the northern
boundary: the indicator in the theory we cite is the freshwater convergence
between the southern and northern boundaries, and $\Fovs$ is a proxy for it
whose adequacy we have not tested. Both are natural next steps and neither is
addressed here.

The second is ECCO-V4r4, and it is a genuine disagreement rather than a caveat.
It is the only product here that is dynamically consistent by construction and
therefore the only one whose freshwater budget closes without assimilation
increments, and it shows no decline in $\Fovs$. Its limb salinities are flat only in the
aggregate: its upper cell salinifies at $+0.022$~PSU per decade ($p=0.025$),
the same sign as the other two products, while a significant redistribution
towards its abyssal cell cancels the effect. Against GLORYS12V1 this is
not an underpowered null: on the window the two share, ECCO's minimum detectable
limb-salinity trend of 0.020~PSU per decade is well below the GLORYS12V1 rate of
$+0.083$, so ECCO could have seen that salinification and did not. Against
ORAS5, whose common-window rate is $+0.019$, the comparison is inconclusive.
ECCO is also the one product in which the coarse limb redistribution is
significant and acts to cancel a positive water-mass term, so its flat limb
salinity is a cancellation rather than an absence of change. The mechanism we report is therefore a
property of the two sequential-assimilation products, and it remains possible
that what they are tracking is the arrival of the Argo array rather than a
change in the ocean. The evidence against that reading is that the attribution
is unchanged in every window and under every inhomogeneity correction we tried,
and that EN4 reproduces the limb asymmetry independently. The evidence for it is
that both sequential-assimilation products step at an observing-system boundary:
GLORYS12V1's decline is a level shift at the 2004 Argo transition rather than a
trend, and the ORAS5 record steps across its consolidated-to-operational join,
significantly so in the limb salinity itself ($+0.038$~PSU, $p=0.029$), after
which its Argo-era trend is no longer significant. Neither product's \emph{rate}
of decline should be taken at face value. We could not obtain the salinity
assimilation increments that would settle the question directly, and we regard
it as open.

Beyond that, all of these products are forced by the ERA family, so
cross-product agreement is a weaker form of independence than the product count
suggests: they differ in ocean model and in assimilation scheme, but not in the
atmosphere driving them. The Copernicus distribution of ORAS5 publishes one of
the system's five ensemble members\cite{Zuo2019}, so cross-product spread has to
stand in for within-product spread throughout, and that spread, 0.094~Sv, exceeds even the
widest within-product interval\cite{Jackson2019}. We compute
ECCO on the distributed 0.5\degree{} interpolation rather than the native LLC90
grid, which does not conserve transport exactly, and since ECCO is now the
load-bearing counterexample that limitation matters more than it would if it
were one supporting vote among four. The three products differ by 46\% in northward-limb transport
at this section (14.6 to 21.3~Sv), and we have not applied a published SAMBA or
SAMOC mean to arbitrate between them. That omission is not neutral: the weakest
of the three is ECCO, which is also our counterexample, so a comparison against
the observed mean transport could as easily undermine the counterexample as
support it, and it should be done before these numbers are used to weigh one
product against another. The reanalyses are also least
constrained by observations exactly where we evaluate them, since the observing
network at 34.5\degS{} is far sparser than the sustained array at
26.5\degN{}\cite{Meinen2018,Cunningham2007,GarzoliMatano2011,FrajkaWilliams2019}.

For models the implication cuts two ways. Correcting the freshwater transport
bias is necessary if a model is to represent the salt-advection feedback with
the right sign, and the CMIP6 median sits above every observational estimate we
computed. But the majority count usually quoted is not distinguishable from a
coin flip once model genealogy is accounted for, and we find no detectable
relationship between the sign of $\Fovs$ and the rate at which a model's
overturning weakens, either over the satellite era or over the forced response
to 2100, with a sample that could not have detected a moderate one. What we do
find is that model $\Fovs$ is strongly correlated with pre-industrial
overturning strength, which suggests that the freshwater transport bias is
partly a symptom of the overturning strength bias. If that is right, correcting
$\Fovs$ in isolation is the wrong target, and the two biases should be
addressed together.

What would settle the transient question is observational rather than
analytical. The claim we make is now a claim about limb salinities, which
sustained hydrography at 34.5\degS{} and the repeat sections that cross it can
test directly, without passing through an assimilating model; projecting an
observation-only salinity analysis onto the same limb mask is the single most
informative next step and we regard it as the test this result should be judged
by. Sustained transport monitoring at this latitude, extended towards the length
RAPID has reached at 26.5\degN{}, would constrain both factors of
equation~\eqref{eq:factorisation} at once. Reanalyses forced outside the ERA
family would test how much of the present agreement comes from shared boundary
conditions, and release of all five ORAS5 ensemble members would separate
assimilation noise from ocean signal within a single product. Until then the
defensible statement is the narrow one: the Atlantic freshwater budget sits
where the salt-advection feedback is self-reinforcing, and in the products that
decline, the recent decline in that diagnostic records a salinifying upper limb
rather than a slowing circulation.
